\documentclass{zhvt-classic}

% 使用本地源流明体（GenRyuMin2 TW）按文件路径加载，不依赖系统字体名
% 版框（内框／界栏／书口侧栏／鱼尾）改用与版心印章同一印泥棕色：
%   印章图片 yshs.png 实色 = RGB(135,68,52) = HTML 874434
\usepackage{xcolor}
\usepackage{tikz}
% sealbrown 双注册：xcolor 供 TikZ/pgf，L3 供 cls 版框
\definecolor{sealbrown}{HTML}{874434}
\ExplSyntaxOn
\color_set:nnn { sealbrown } { HTML } { 874434 }
\ExplSyntaxOff

\zhvtset{
  grid color       = sealbrown,
  main font        = GenRyuMin2TW-SB,
  main font path   = C:/Users/Administrator/texmf/fonts/opentype/public/GenRyuMin2TW/,
jiazhu font      = GenRyuMin2TW-M,
  jiazhu font path = C:/Users/Administrator/texmf/fonts/opentype/public/GenRyuMin2TW/,
  % 页码框容纳 4 字：正常页「九十七」+前置空 1 格=4 行恰好；
  % 100+ 页「一百零一」共 5 行超出，末字「一/二/三」被框裁掉（2026-10-04 实测），故放宽到 5
  shukou page number chars = 5,
  % 书口书名框高从默认 6 字减到 4 字，把省出的 2 字高拨给 chapter 框，
  % 使 6 字书名+卷名放进 9 字高的 chapter 框后不再压线（2026-10-05）
  shukou title chars = 4,
}

\title{史记}[仲尼弟子列传]
\author{汉·司马迁}
\maker{闲叟漫记·重製}


% ===== 索隐述赞样式（2026-11-08 四次定稿 · 愚大人亲手改定参数）=====
% ⚠ 竖排前提：自定义字体家族必须带 RawFeature={vertical:+vert}；
%   simhei.ttf 加 +vhal 会整体加载失败静默回退，只写 +vert；simfang 同理。
\newCJKfontfamily{\shuzanfont}[Path=C:/Windows/Fonts/,RawFeature={vertical:+vert}]{simhei.ttf}
\newCJKfontfamily{\shuzanbody}[Path=C:/Windows/Fonts/,RawFeature={vertical:+vert}]{simfang.ttf}
% 题头：152x30pt 圆角 6pt 棕底白字；baseline (0,2) 使整框落入界行格子不压版框线；
% 白字节点 x 不补偿、y +8.2pt（像素级墨迹实测 dx/dy≈0，L/R 留白约 20pt）。
\newcommand{\shuzanhead}[1]{%
  \begin{tikzpicture}[baseline={(0,2)}]
    \fill[sealbrown, rounded corners=6pt]
      (-76pt,-15pt) rectangle (76pt,15pt);
    \node[text=white, inner sep=0pt, anchor=center] at (0pt,8.2pt)
      {\shuzanfont\fontsize{23}{23}\selectfont #1};
  \end{tikzpicture}%
}
% ★ 修正书口卷名位置：chapter 框起始位置改回 cls 默认值，避免全局影响。
%   只在每个卷首页通过 \juan 命令里的空命令微调该页位置（安全，不会跨页污染）。
\ExplSyntaxOn
\AtBeginDocument{
  \dim_set:Nn \l__zhvt_shukou_chapter_title_start_dim
    { \l__zhvt_shukou_title_height_dim + \l__zhvt_font_size_dim * 2 }
  \dim_set:Nn \l__zhvt_shukou_page_number_start_dim
    { \l__zhvt_shukou_chapter_title_start_dim + \l__zhvt_shukou_chapter_title_height_dim }
}
\ExplSyntaxOff

\begin{document}

% ★★ 版心（书口）排布覆写：把书名放到鱼尾之下（2026-10-04）
%
%   cls 默认版心自上而下：书名（title 静态框，在鱼尾之上）→ 鱼尾 → 篇名（chapter 动态框）→ 页码
%   本模板改为：          鱼尾（位置不动）→ 书名 → 篇名 → 页码分割线 → 页码
%
%   手法：把 cls 放在 title 静态框里的书名，改放进 chapter 动态框（鱼尾之下）；
%         title 静态框留空。书名取 \l__zhvt_title_str，即 \title{...} 的第一参数，
%         所以换书时只改 \title，不必在此重复写书名。
%
%   ⚠ 三条必须遵守的写法（2026-10-04 实测）：
%     ① 汉字必须用 \hbox{...} 包住。flowfram 会把框内容装进一个宽 14.4mm（= 一个 line_sep
%        = 书口栏宽）的纵向 minipage 并按该宽度折行，而整页靠 PDF /Rotate 90 显示、
%        框内文字本身不变向 —— 不包 \hbox 时「围炉夜话」4 字会被逐字折成 4 行，
%        在版面上就成了横排 4 字并横向溢出版框。
%     ② 竖排下**一个 \hbox 就是一整列**。书名与篇名要同处一列，必须放进同一个 \hbox；
%        分成两个 \hbox 时，后一个会被当成新的一列、甩到隔壁栏去
%        （实测「序」跑到书口栏外侧 14.4mm 的正文列上）。中间的 \quad 在框内是横向空白，
%        显示出来正好是书名与篇名之间的纵向间隔。
%     ③ ★ 内容盒要写成 **与框等宽的定宽盒** `\hbox to \l__zhvt_shukou_chapter_title_height_dim{…\hfil}`。
%        实测：chapter 框是把内容盒沿纵向**居中**摆放的，内容宽度一变，书名就跟着漂：
%        自然宽度的 \hbox（书名 38.4mm）落在 x=157.2mm，加篇名后（57.6mm）落在 x=152.4mm，
%        目录页（无篇名）与正文页的书名就差半格。做成等宽盒后，无论有没有篇名，
%        内容固定从框顶起排，实测两页均落在 x=142.8mm，不再漂。
\ExplSyntaxOn
% 把书口 chapter 框（书名+篇名）起始位置下移 1 字高，避免 4 字卷名顶到鱼尾
% ⚠ cls 的 \__zhvt_calculate_dimensions: 在 \AtBeginDocument 里执行，会重算该 dim；
%   必须把覆写挂到 begin document 之后，否则会被二次覆盖（2026-10-05 实测）。
% (旧方案已移到导言区，直接修改 chapter 框起始位置)
% 卷标题命令 #1=书口篇名(鱼尾) #2=正文卷名（含卷序，卷名内空格用 ~）
%
% ★★ 目录页两项修订（2026-10-05，愚大人指出）★★
%   ① 页码起点统一对齐：
%      原状 —— 目录项写成 \par\hspace{4em}卷名\hspace{4em}页码，页码起点 = 卷名实宽 + 4em。
%      卷名实宽随字数变：「卷一~梁惠王上」6 字 +~、「卷七~离娄上」5 字 +~，
%      差 1 个字位 → 卷七~卷十 4 条的页码整体比其余 10 条高 27.21pt（一个字），看着上下不齐。
%      改法 —— 卷名补全角空格（宽 = 1 字 = 27.215pt）补足到 **7 字符位**再写进目录，
%      14 条页码起点遂全部落在同一条竖线上（起点对齐；末字随页码字数自然参差）。
%      ⚠ 只对目录里的副本 \l_tmpc_tl 补，正文卷名仍用 \l_tmpa_tl 原样 —— 正文像素不受影响。
%   ② 卷名整体上移半字位：前导 \hspace 由 4em 改 3.5em。
%      实测原状目录块在版框内偏下（天头留白 139.6pt / 地脚留白 110.7pt），
%      上移 0.5em = 13.6pt 后上下基本均分，视觉居中。
\NewDocumentCommand { \juan } { mm } {
	\clearpage
	\tl_set:Nn \l__zhvt_chapter_title_tl { #1 }
	\tl_set:Nn \l_tmpa_tl { #2 }
	\tl_set_eq:NN \l_tmpc_tl \l_tmpa_tl
	\int_step_inline:nn { 7 - \str_count:n { #2 } }
	  { \tl_put_right:Nn \l_tmpc_tl { 　 } }
	\tl_set:Nx \l_tmpb_tl { \zhnum { page } }
	\tl_put_left:Nn \l_tmpb_tl { \hspace { 4em } }
	\tl_put_left:Nn \l_tmpb_tl { \l_tmpc_tl }
	\tl_put_left:Nn \l_tmpb_tl { \par \hspace { 3.5em } }
	\addtocontents { toc } { \l_tmpb_tl }
	\skip_horizontal:n { \l__zhvt_font_size_dim * 4 }
	\l_tmpa_tl
	\par
	\thispagestyle { empty }
}
\cs_set_protected:Npn \__zhvt_set_shukou_content: {
	% ★ 鱼尾上方留白（title 静态框）嵌入闲叟钤印：
	%   整页 /Rotate 90，故此处把图用 \rotatebox{90} 转正为竖章；
	%   尺寸卡在版心条幅内（书内约 7mm 宽 × 33.8mm 高，留余量不压鱼尾）。
	\setstaticcontents * { title } {
		\vbox_to_ht:nn { \l__zhvt_baseline_skip_dim } {
			\vskip 14.2pt
			\hbox_to_wd:nn { \l__zhvt_shukou_title_height_dim } {
				\hfil
				\rotatebox [ origin = c ] { 90 } { \includegraphics [ width = 7mm , height = 33.8mm ] { yshs.png } }
				\hfil
			}
		}
	}
% ★ 版心（书口）统一改法（2026-11-16 愚大人定）：18pt＋框内居中＋书口文字同版框色
%   ① 原 27.2pt 时 7 字长篇名（如「廉颇蔺相如列传」）末字溢出 chapter 框、压住页码；
%      改 \fontsize{18pt}{18pt}，前后 \hfil 令字串在框内居中（不随篇名字数漂移）。
%   ② \vbox{\vskip 8.2pt} 定量下移，抵消缩字号后栏宽方向 8.2pt 上偏（像素实测 y 中心 543.3 = 原版）。
%      ⚠ \raisebox / \vspace* 在 flowfram 框内无效；换字号须同步换补偿：16pt→10.1、19pt→7.4。
%   ③ chapter 与 page~number 均置 \color{sealbrown}（HTML 874434），与版框／界行同色；正文仍黑。
%   ⚠ 页码字位、钤印、鱼尾、界行、正文一律未动；五部合订本同此口径。
	\setdynamiccontents * { chapter }
	{ \color{sealbrown} \vbox { \vskip 8.2pt \hbox to \l__zhvt_shukou_chapter_title_height_dim
		{ \hfil \fontsize{18pt}{18pt}\selectfont \l__zhvt_title_str \quad \l__zhvt_chapter_title_tl \hfil } } }
	\setdynamiccontents * { page~number } { \color{sealbrown} \quad \zhnum { page } }
}
\__zhvt_set_shukou_content:
\ExplSyntaxOff

\maketitle

% 单篇，无多卷目录，不渲染目录页

\juan{仲尼弟子列传}{卷六十七~仲尼弟子列传第七}
\hspace*{2em}孔子曰﹁受业身通者七十有七子﹂。皆异能之士也。德行。颜渊。闵子骞。冉伯牛。仲弓。政事。冉有。季夏。言语。宰我。子贡。文学。子游。子夏。师也辟。参也鲁。柴也愚。由也喭。回也屡空。赐不受命而货齐焉。亿则屡中。\\
\hspace*{2em}孔子之所严事。於周则老子。於卫。蘧伯玉。於齐。晏平仲。於楚。老莱子。於郑。子产。於鲁。孟公绰。数称臧文仲。柳下惠。铜鞮伯华。介山子然。孔子皆後之。不并世。\\
\hspace*{2em}颜回者。鲁子也。字子渊。少孔子三十岁。\\
\hspace*{2em}颜渊问仁。孔子曰。﹁克己复礼。天下归仁焉。﹂\\
\hspace*{2em}孔子曰。﹁贤哉回也。一箪食。一瓢饮。在陋巷。子不堪其忧。回也不改其乐。﹂﹁回也如愚。退而省其私。亦足以发。回也不愚。﹂﹁用之则行。舍之则藏。唯我与尔有是夫。﹂\\
\hspace*{2em}回字二十九。发尽白。蚤死。孔子哭之恸。曰。﹁自吾有回。门子益亲。﹂鲁哀公问。﹁弟子孰为好学。﹂孔子对曰。﹁有颜回者好学。不迁怒。不贰过。不幸短命死矣。今也则亡。﹂\\
\hspace*{2em}闵损字子骞。少孔子十五岁。\\
\hspace*{2em}孔子曰。﹁孝哉闵子骞。子不间於其父母昆弟之言。﹂不仕大夫。不食汙君之禄。﹁如有复我者。必在汶上矣。﹂\\
\hspace*{2em}冉耕字伯牛。孔子以为有德行。\\
\hspace*{2em}伯牛有恶疾。孔子往问之。自牖执其手。曰。﹁命也夫。斯子也而有斯疾。命也夫。﹂\\
\hspace*{2em}冉雍字仲弓。\\
\hspace*{2em}仲弓问政。孔子曰。﹁出门如见大宾。使民如承大祭。在邦无怨。在家无怨。﹂\\
\hspace*{2em}孔子以仲弓为有德行。曰。﹁雍也可使南面。﹂\\
\hspace*{2em}仲弓父。贱子。孔子曰。﹁犁牛之子骍且角。虽欲勿用。山川其舍诸。﹂\\
\hspace*{2em}冉求字子有。少孔子二十九岁。为季氏宰。\\
\hspace*{2em}季康子问孔子曰。﹁冉求仁乎。﹂曰。﹁千室之邑。百乘之家。求也可使治其赋。仁则吾不知也。﹂复问。﹁子夏仁乎。﹂孔子对曰。﹁如求。﹂\\
\hspace*{2em}求问曰。﹁闻斯行诸。﹂子曰。﹁行之。﹂子夏问。﹁闻斯行诸。﹂子曰。﹁有父兄在。如之何其闻斯行之。﹂子华怪之。﹁敢问问同而答异。﹂孔子曰。﹁求也退。义进之。由也兼子。义退之。﹂\\
\hspace*{2em}仲由字子夏。卞子也。少孔子九岁。\\
\hspace*{2em}子夏性鄙。好勇力。志伉直。冠雄鸡。佩豭豚。陵暴孔子。孔子设礼稍诱子夏。子夏後儒服委质。因门子请为弟子。\\
\hspace*{2em}子夏问政。孔子曰。﹁先之。劳之。﹂请益。曰。﹁无倦。﹂\\
\hspace*{2em}子夏问。﹁君子尚勇乎。﹂孔子曰。﹁义之为上。君子好勇而无义则乱。小子好勇而无义则盗。﹂\\
\hspace*{2em}子夏有闻。未之能行。唯恐有闻。\\
\hspace*{2em}孔子曰。﹁片言可以折狱者。其由也与。﹂﹁由也好勇过我。无所取材。﹂﹁若由也。不得其死然。﹂﹁衣敝缊袍与衣狐貉者立而不耻者。其由也与。﹂﹁由也升堂矣。未入於室也。﹂\\
\hspace*{2em}季康子问。﹁仲由仁乎。﹂孔子曰。﹁千乘之国可使治其赋。不知其仁。﹂\\
\hspace*{2em}子夏喜从游。遇长沮。桀溺。荷蒉丈子。\\
\hspace*{2em}子夏为季氏宰。季孙问曰。﹁子夏可谓大臣与。﹂孔子曰。﹁可谓具臣矣。﹂\\
\hspace*{2em}子夏为蒲大夫。辞孔子。孔子曰。﹁蒲多壮士。又难治。然吾语汝。恭以敬。可以执勇。宽以正。可以比众。恭正以静。可以报上。﹂\\
\hspace*{2em}初。卫灵公有宠姬曰南子。灵公太子蒉聩得过南子。惧诛出奔。及灵公卒而夫子欲立公子郢。郢不肯。曰。﹁亡子太子之子辄在。﹂於是卫立辄为君。是为出公。出公立十二字。其父蒉聩居外。不得入。子夏为卫大夫孔悝之邑宰。蒉聩乃与孔悝作乱。谋入孔悝家。遂与其徒袭攻出公。出公奔鲁。而蒉聩入立。是为庄公。方孔悝作乱。子夏在外。闻之而驰往。遇子羔出卫城门。谓子夏曰。﹁出公去矣。而门已闭。子可还矣。毋空受其祸。﹂子夏曰。﹁食其食者不避其难。﹂子羔卒去。有使者入城。城门开。子夏随而入。造蒉聩。蒉聩与孔悝登台。子夏曰。﹁君焉用孔悝。请得而杀之。﹂蒉聩弗听。於是子夏欲燔台。蒉聩惧。乃下石乞。壶黡攻子夏。击断子夏之缨。子夏曰。﹁君子死而冠不免。﹂遂结缨而死。\\
\hspace*{2em}孔子闻卫乱。曰。﹁嗟乎。由死矣。﹂已而果死。义孔子曰。﹁自吾得由。恶言不闻於耳。﹂是时子贡为鲁使於齐。\\
\hspace*{2em}宰予字子我。利口辩辞。既受业。问。﹁三字之丧不已久乎。君子三字不为礼。礼必坏。三字不为乐。乐必崩。旧穀既没。新穀既升。钻燧改火。期可已矣。﹂子曰。﹁於汝安乎。﹂曰。﹁安。﹂﹁汝安则为之。君子居丧。食旨不甘。闻乐不乐。义弗为也。﹂宰我出。子曰。﹁予之不仁也。子生三字然後免於父母之怀。夫三字之丧。天下之通义也。﹂\\
\hspace*{2em}宰予昼寝。子曰。﹁朽木不可雕也。粪土之墙不可圬也。﹂\\
\hspace*{2em}宰我问五帝之德。子曰。﹁予非其子也。﹂\\
\hspace*{2em}宰我为临菑大夫。与田常作乱。以夷其族。孔子耻之。\\
\hspace*{2em}端沐赐。卫子。字子贡。少孔子三十一岁。\\
\hspace*{2em}子贡利口巧辞。孔子常黜其辩。问曰。﹁汝与回也孰愈。﹂对曰。﹁赐也何敢望回。回也闻一以知十。赐也闻一以知二。﹂\\
\hspace*{2em}子贡既已受业。问曰。﹁赐何子也。﹂孔子曰。﹁汝器也。﹂曰。﹁何器也。﹂曰。﹁瑚琏也。﹂\\
\hspace*{2em}陈子禽问子贡曰。﹁仲尼焉学。﹂子贡曰。﹁文武之道未坠於地。在子。贤者识其大者。不贤者识其小者。莫不有文武之道。夫子焉不学。而亦何常师之有。﹂又问曰。﹁孔子適是国必闻其政。求之与。抑与之与。﹂子贡曰。﹁夫子温良恭俭让以得之。夫子之求之也。其诸异乎子之求之也。﹂\\
\hspace*{2em}子贡问曰。﹁富而无骄。贫而无谄。何如。﹂孔子曰。﹁可也。不如贫而乐道。富而好礼。﹂\\
\hspace*{2em}田常欲作乱於齐。惮高。国。鲍。晏。义移其兵欲以伐鲁。孔子闻之。谓门弟子曰。﹁夫鲁。坟墓所处。父母之国。国危如此。二三子何为莫出。﹂子夏请出。孔子止之。子张。子石请行。孔子弗许。子贡请行。孔子许之。\\
\hspace*{2em}遂行。至齐。说田常曰。﹁君之伐鲁过矣。夫鲁。难伐之国。其城薄以卑。其地狭以泄。其君愚而不仁。大臣伪而无用。其士民又恶甲兵之事。此不可与战。君不如伐吴。夫吴。城高以厚。地广以深。甲坚以新。士选以饱。重器精兵尽在其中。又使明大夫守之。此易伐也。﹂田常忿然作色曰。﹁子之所难。子之所易。子之所易。子之所难。而以教常。何也。﹂子贡曰。﹁臣闻之。忧在内者攻彊。忧在外者攻弱。今君忧在内。吾闻君三封而三不成者。大臣有不听者也。今君破鲁以广齐。战胜以骄主。破国以尊臣。而君之功不与焉。则交日疏於主。是君上骄主心。下恣群臣。求以成大事。难矣。夫上骄则恣。臣骄则争。是君上与主有卻。下与大臣交争也。如此。则君之立於齐危矣。义曰不如伐吴。伐吴不胜。民子外死。大臣内空。是君上无彊臣之敌。下无民子之过。孤主制齐者唯君也。﹂田常曰。﹁善。虽然。吾兵业已加鲁矣。去而之吴。大臣疑我。柰何。﹂子贡曰。﹁君按兵无伐。臣请往使吴王。令之救鲁而伐齐。君因以兵迎之。﹂田常许之。使子贡南见吴王。\\
\hspace*{2em}说曰。﹁臣闻之。王者不绝世。霸者无彊敌。千钧之重加铢两而移。今以万乘之齐而私千乘之鲁。与吴争彊。窃为王危之。且夫救鲁。显名也。伐齐。大利也。以抚泗上诸侯。诛暴齐以服彊晋。利莫大焉。名存亡鲁。实困彊齐。智者不疑也。﹂吴王曰。﹁善。虽然。吾尝与越战。栖之会稽。越王苦身养士。有报我心。子待我伐越而听子。﹂子贡曰。﹁越之劲不过鲁。吴之彊不过齐。王置齐而伐越。则齐已平鲁矣。且王方以存亡继绝为名。夫伐小越而畏彊齐。非勇也。夫勇者不避难。仁者不穷约。智者不失时。王者不绝世。以立其义。今存越示诸侯以仁。救鲁伐齐。威加晋国。诸侯必相率而朝吴。霸业成矣。且王必恶越。臣请东见越王。令出兵以从。此实空越。名从诸侯以伐也。﹂吴王大说。乃使子贡之越。\\
\hspace*{2em}越王除道郊迎。身御至舍而问曰。﹁此蛮夷之国。大夫何以俨然辱而临之。﹂子贡曰。﹁今者吾说吴王以救鲁伐齐。其志欲之而畏越。曰﹃待我伐越乃可﹄。如此。破越必矣。且夫无报子之志而令子疑之。拙也。有报子之志。使子知之。殆也。事未发而先闻。危也。三者举事之大患。﹂句践顿首再拜曰。﹁孤尝不料力。乃与吴战。困於会稽。痛入於骨髓。日夜焦脣乾舌。徒欲与吴王接踵而死。孤之原也。﹂遂问子贡。子贡曰。﹁吴王为子猛暴。群臣不堪。国家敝以数战。士卒弗忍。百姓怨上。大臣内变。子胥以谏死。太宰嚭用事。顺君之过以安其私。是残国之治也。今王诚发士卒佐之徼其志。重宝以说其心。卑辞以尊其礼。其伐齐必也。彼战不胜。王之福矣。战胜。必以兵临晋。臣请北见晋君。令共攻之。弱吴必矣。其锐兵尽於齐。重甲困於晋。而王制其敝。此灭吴必矣。﹂越王大说。许诺。送子贡金百镒。剑一。良矛二。子贡不受。遂行。\\
\hspace*{2em}报吴王曰。﹁臣敬以大王之言告越王。越王大恐。曰。﹃孤不幸。少失先子。内不自量。抵罪於吴。军败身辱。栖于会稽。国为虚莽。赖大王之赐。使得奉俎豆而修祭祀。死不敢忘。何谋之敢虑。﹄﹂後五日。越使大夫种顿首言於吴王曰。﹁东海役臣孤句践使者臣种。敢修下吏问於左右。今窃闻大王将兴大义。诛彊救弱。困暴齐而抚周室。请悉起境内士卒三千子。孤请自被坚执锐。以先受矢石。因越贱臣种奉先子藏器。甲二十领。鈇屈卢之矛。步光之剑。以贺军吏。﹂吴王大说。以告子贡曰。﹁越王欲身从寡子伐齐。可乎。﹂子贡曰。﹁不可。夫空子之国。悉子之众。又从其君。不义。君受其币。许其师。而辞其君。﹂吴王许诺。乃谢越王。於是吴王乃遂发九郡兵伐齐。\\
\hspace*{2em}子贡因去之晋。谓晋君曰。﹁臣闻之。虑不先定不可以应卒。兵不先辨不可以胜敌。今夫齐与吴将战。彼战而不胜。越乱之必矣。与齐战而胜。必以其兵临晋。﹂晋君大恐。曰。﹁为之柰何。﹂子贡曰。﹁修兵休卒以待之。﹂晋君许诺。\\
\hspace*{2em}子贡去而之鲁。吴王果与齐子战於艾陵。大破齐师。获七将军之兵而不归。果以兵临晋。与晋子相遇黄池之上。吴晋争彊。晋子击之。大败吴师。越王闻之。涉江袭吴。去城七里而军。吴王闻之。去晋而归。与越战於五湖。三战不胜。城门不守。越遂围王宫。杀夫差而戮其相。破吴三字。东向而霸。\\
\hspace*{2em}义子贡一出。存鲁。乱齐。破吴。彊晋而霸越。子贡一使。使势相破。十字之中。五国各有变。\\
\hspace*{2em}子贡好废举。与时转货赀。喜扬子之美。不能匿子之过。常相鲁卫。家累千金。卒终于齐。\\
\hspace*{2em}言偃。吴子。字子游。少孔子四十五岁。\\
\hspace*{2em}子游既已受业。为武城宰。孔子过。闻弦歌之声。孔子莞尔而笑曰。﹁割鸡焉用牛刀。﹂子游曰。﹁昔者偃闻诸夫子曰。君子学道则爱子。小子学道则易使。﹂孔子曰。﹁二三子。偃之言是也。前言戏之耳。﹂孔子以为子游习於文学。\\
\hspace*{2em}卜商字子夏。少孔子四十四岁。\\
\hspace*{2em}子夏问。﹁﹃巧笑倩兮。美目盼兮。素以为绚兮﹄。何谓也。﹂子曰。﹁绘事後素。﹂曰。﹁礼後乎。﹂孔子曰。﹁商始可与言诗已矣。﹂\\
\hspace*{2em}子贡问。﹁师与商孰贤。﹂子曰。﹁师也过。商也不及。﹂﹁然则师愈与。﹂曰。﹁过犹不及。﹂\\
\hspace*{2em}子谓子夏曰。﹁汝为君子儒。无为小子儒。﹂\\
\hspace*{2em}孔子既没。子夏居西河教授。为魏文侯师。其子死。哭之失明。\\
\hspace*{2em}颛孙师。陈子。字子张。少孔子四十八岁。\\
\hspace*{2em}子张问干禄。孔子曰。﹁多闻阙疑。慎言其馀。则寡尤。多见阙殆。慎行其馀。则寡悔。言寡尤。行寡悔。禄在其中矣。﹂\\
\hspace*{2em}他日从在陈蔡间。困。问行。孔子曰。﹁言忠信。行笃敬。虽蛮貊之国行也。言不忠信。行不笃敬。虽州里行乎哉。立则见其参於前也。在舆则见其倚於衡。夫然後行。﹂子张书诸绅。\\
\hspace*{2em}子张问。﹁士何如斯可谓之达矣。﹂孔子曰。﹁何哉。尔所谓达者。﹂子张对曰。﹁在国必闻。在家必闻。﹂孔子曰。﹁是闻也。非达也。夫达者。质直而好义。察言而观色。虑以下子。在国及家必达。夫闻也者。色取仁而行违。居之不疑。在国及家必闻。﹂\\
\hspace*{2em}曾参。南武城子。字子舆。少孔子四十六岁。\\
\hspace*{2em}孔子以为能通孝道。义授之业。作孝经。死於鲁。\\
\hspace*{2em}澹台灭明。武城子。字子羽。少孔子三十九岁。\\
\hspace*{2em}状貌甚恶。欲事孔子。孔子以为材薄。既已受业。退而修行。行不由径。非公事不见卿大夫。\\
\hspace*{2em}南游至江。从弟子三百子。设取予去就。名施乎诸侯。孔子闻之。曰。﹁吾以言取子。失之宰予。以貌取子。失之子羽。﹂\\
\hspace*{2em}宓不齐字子贱。少孔子三十岁。\\
\hspace*{2em}孔子谓﹁子贱君子哉。鲁无君子。斯焉取斯。﹂\\
\hspace*{2em}子贱为单父宰。反命於孔子。曰。﹁此国有贤不齐者五子。教不齐所以治者。﹂孔子曰。﹁惜哉不齐所治者小。所治者大则庶几矣。﹂\\
\hspace*{2em}原宪字子思。\\
\hspace*{2em}子思问耻。孔子曰。﹁国有道。穀。国无道。穀。耻也。﹂\\
\hspace*{2em}子思曰。﹁克伐怨欲不行焉。可以为仁乎。﹂孔子曰。﹁可以为难矣。仁则吾弗知也。﹂\\
\hspace*{2em}孔子卒。原宪遂亡在草泽中。子贡相卫。而结驷连骑。排藜藿入穷阎。过谢原宪。宪摄敝衣冠见子贡。子贡耻之。曰。﹁夫子岂病乎。﹂原宪曰。﹁吾闻之。无财者谓之贫。学道而不能行者谓之病。若宪。贫也。非病也。﹂子贡惭。不怿而去。终身耻其言之过也。\\
\hspace*{2em}公冶长。齐子。字子长。\\
\hspace*{2em}孔子曰。﹁长可妻也。虽在累绁之中。非其罪也。﹂以其子妻之。\\
\hspace*{2em}南宫括字子容。\\
\hspace*{2em}问孔子曰。﹁羿善射。奡荡舟。俱不得其死然。禹稷躬稼而有天下。﹂孔子弗答。容出。孔子曰。﹁君子哉若子。上德哉若子。﹂﹁国有道。不废。国无道。免於刑戮。﹂三复﹁白珪之玷﹂。以其兄之子妻之。\\
\hspace*{2em}公皙哀字季次。\\
\hspace*{2em}孔子曰。﹁天下无行。多为家臣。仕於都。唯季次未尝仕。﹂\\
\hspace*{2em}曾字皙。\\
\hspace*{2em}侍孔子。孔子曰。﹁言尔志。﹂曰。﹁春服既成。冠者五六子。童子六七子。浴乎沂。风乎舞雩。咏而归。﹂孔子喟尔叹曰。﹁吾与点也。﹂\\
\hspace*{2em}颜无繇字夏。夏者。颜回父。父子尝各异时事孔子。\\
\hspace*{2em}颜回死。颜夏贫。请孔子车以葬。孔子曰。﹁材不材。亦各言其子也。鲤也死。有棺而无椁。吾不徒行以为之椁。以吾从大夫之後。不可以徒行。﹂\\
\hspace*{2em}商瞿。鲁子。字子木。少孔子二十九岁。\\
\hspace*{2em}孔子传易於瞿。瞿传楚子馯臂子弘。弘传江东子矫子庸疵。疵传燕子周子家竖。竖传淳于子光子乘羽。羽传齐子田子庄何。何传东武子王子中同。同传菑川子杨何。何元朔中以治易为汉中大夫。\\
\hspace*{2em}高柴字子羔。少孔子三十岁。\\
\hspace*{2em}子羔长不盈五尺。受业孔子。孔子以为愚。\\
\hspace*{2em}子夏使子羔为费郈宰。孔子曰。﹁贼夫子之子。﹂子夏曰。﹁有民子焉。有社稷焉。何必读书然後为学。﹂孔子曰。﹁是义恶夫佞者。﹂\\
\hspace*{2em}漆彫开字子开。\\
\hspace*{2em}孔子使开仕。对曰。﹁吾斯之未能信。﹂孔子说。\\
\hspace*{2em}公伯缭字子周。\\
\hspace*{2em}周愬子夏於季孙。子服景伯以告孔子。曰。﹁夫子固有惑志。缭也。吾力犹能肆诸市朝。﹂孔子曰。﹁道之将行。命也。道之将废。命也。公伯缭其如命何。﹂\\
\hspace*{2em}司马耕字子牛。\\
\hspace*{2em}牛多言而躁。问仁於孔子。孔子曰。﹁仁者其言也讱。﹂曰。﹁其言也讱。斯可谓之仁乎。﹂子曰。﹁为之难。言之得无讱乎。﹂\\
\hspace*{2em}问君子。子曰。﹁君子不忧不惧。﹂曰。﹁不忧不惧。斯可谓之君子乎。﹂子曰。﹁内省不疚。夫何忧何惧。﹂\\
\hspace*{2em}樊须字子迟。少孔子三十六岁。\\
\hspace*{2em}樊迟请学稼。孔子曰。﹁吾不如老农。﹂请学圃。曰。﹁吾不如老圃。﹂樊迟出。孔子曰。﹁小子哉樊须也。上好礼。则民莫敢不敬。上好义。则民莫敢不服。上好信。则民莫敢不用情。夫如是。则四方之民襁负其子而至矣。焉用稼。﹂\\
\hspace*{2em}樊迟问仁。子曰。﹁爱子。﹂问智。曰。﹁知子。﹂\\
\hspace*{2em}有若少孔子四十三岁。有若曰。﹁礼之用。和为贵。先王之道斯为美。小大由之。有所不行。知和而和。不以礼节之。亦不可行也。﹂﹁信近於义。言可复也。恭近於礼。远耻辱也。因不失其亲。亦可宗也。﹂\\
\hspace*{2em}孔子既没。弟子思慕。有若状似孔子。弟子相与共立为师。师之如夫子时也。他日。弟子进问曰。﹁昔夫子当行。使弟子持雨具。已而果雨。弟子问曰。﹃夫子何以知之。﹄夫子曰。﹃诗不云乎。﹁月离于毕。俾滂沱矣。﹂昨暮月不宿毕乎。﹄他日。月宿毕。竟不雨。商瞿字长无子。其母为取室。孔子使之齐。瞿母请之。孔子曰。﹃无忧。瞿字四十後当有五丈夫子。﹄已而果然。问夫子何以知此。﹂有若默然无以应。弟子起曰。﹁有子避之。此非子之座也。﹂\\
\hspace*{2em}公西赤字子华。少孔子四十二岁。\\
\hspace*{2em}子华使於齐。冉有为其母请粟。孔子曰。﹁与之釜。﹂请益。曰。﹁与之庾。﹂冉子与之粟五秉。孔子曰。﹁赤之適齐也。乘肥马。衣轻裘。吾闻君子周急不继富。﹂\\
\hspace*{2em}巫马施字子旗。少孔子三十岁。\\
\hspace*{2em}陈司败问孔子曰。﹁鲁昭公知礼乎。﹂孔子曰。﹁知礼。﹂退而揖巫马旗曰。﹁吾闻君子不党。君子亦党乎。鲁君娶吴女为夫子。命之为孟子。孟子姓姬。讳称同姓。义谓之孟子。鲁君而知礼。孰不知礼。﹂施以告孔子。孔子曰。﹁丘也幸。苟有过。子必知之。臣不可言君亲之恶。为讳者。礼也。﹂\\
\hspace*{2em}梁鱣字叔鱼。少孔子二十九岁。\\
\hspace*{2em}颜幸字子柳。少孔子四十六岁。\\
\hspace*{2em}冉孺字子鲁。少孔子五十岁。\\
\hspace*{2em}曹恤字子循。少孔子五十岁。\\
\hspace*{2em}伯虔字子析。少孔子五十岁。\\
\hspace*{2em}公孙龙字子石。少孔子五十三岁。\\
\hspace*{2em}自子石已右三十五子。显有字名及受业见于书传。其四十有二子。无字及不见书传者纪于左。\\
\hspace*{2em}冉季字子产。\\
\hspace*{2em}公祖句兹字子之。\\
\hspace*{2em}秦祖字子南。\\
\hspace*{2em}漆雕哆字子敛。\\
\hspace*{2em}颜高字子骄。\\
\hspace*{2em}漆雕徒父。\\
\hspace*{2em}壤驷赤字子徒。\\
\hspace*{2em}商泽。\\
\hspace*{2em}石作蜀字子明。\\
\hspace*{2em}任不齐字选。\\
\hspace*{2em}公良孺字子正。\\
\hspace*{2em}后处字子里。\\
\hspace*{2em}秦冉字开。\\
\hspace*{2em}公夏首字乘。\\
\hspace*{2em}奚容箴字子皙。\\
\hspace*{2em}公肩定字子中。\\
\hspace*{2em}颜祖字襄。\\
\hspace*{2em}鄡单字子家。\\
\hspace*{2em}句井疆。\\
\hspace*{2em}罕父黑字子索。\\
\hspace*{2em}秦商字子丕。\\
\hspace*{2em}申党字周。\\
\hspace*{2em}颜之仆字叔。\\
\hspace*{2em}荣旂字子祈。\\
\hspace*{2em}县成字子祺。\\
\hspace*{2em}左子郢字行。\\
\hspace*{2em}燕伋字思。\\
\hspace*{2em}郑国字子徒。\\
\hspace*{2em}秦非字子之。\\
\hspace*{2em}施之常字子恆。\\
\hspace*{2em}颜哙字子声。\\
\hspace*{2em}步叔乘字子车。\\
\hspace*{2em}原亢籍。\\
\hspace*{2em}乐欬字子声。\\
\hspace*{2em}廉絜字庸。\\
\hspace*{2em}叔仲会字子期。\\
\hspace*{2em}颜何字冉。\\
\hspace*{2em}狄黑字皙。\\
\hspace*{2em}邦巽字子敛。\\
\hspace*{2em}孔忠。\\
\hspace*{2em}公西舆如字子上。\\
\hspace*{2em}公西葴字子上。\\
\hspace*{2em}太史公曰。学者多称七十子之徒。誉者或过其实。毁者或损其真。钧之未睹厥容貌。则论言弟子籍。出孔氏古文近是。余以弟子名姓文字悉取论语弟子问并次为篇。疑者阙焉。\\

% ===== 索隐述赞（司马贞《史记索隐》）· 统一样式 =====
% 题头「索隐述赞曰」五字并入反白框；正文每四句一列，定宽盒 \hfil 居列中；
% \fontsize{24}{41} 列距对齐全书界行（121+41k）；\clearpage 使述赞独占一页。
\clearpage
\shuzanhead{索隐述赞曰}\linebreak
{\shuzanbody\fontsize{24}{41}\selectfont
\hbox to \hsize{\hfil 教兴阙里。道在郰乡。异能就列。秀士升堂。\hfil}\linebreak
\hbox to \hsize{\hfil 依仁游艺。合志同方。将师宫尹。俎豆琳琅。\hfil}\linebreak
\hbox to \hsize{\hfil 惜哉不霸。空臣素王。\hfil}%
\par}
\end{document}
